\documentclass{article}
\usepackage{graphicx} % Required for inserting images
\usepackage{amsmath}
\usepackage{amssymb}
\usepackage{parskip} % remove global indent
\usepackage{float} % Add to your preamble

\begin{document}


\section*{Question 1}

% --------------------------------- QUESTION 
\subsection*{1.a}

The accepted string in $L$ should be simplified first. Observe that since $w=(aba)^x$ for any $x \geq 0$, we can rewrite $w^R$ as:

$$w^R \equiv ((aba)^x)^R \equiv ((aba)^R)^x) \equiv (aba)^x \equiv w$$

And $L$ can be modified to:

$$L = \{www | w \in (aba)^*\}$$
$$L = \{w^3 | w \in (aba)^*\}$$

Assume $L$ is regular, then there exists a pumping length $k$ for $L$. Let's choose $s = (aba)^{3k}$, $|s| \geq k$; $s \in L$ since $(aba)^{3k} \equiv ((aba)^{k})^3 \in L$.

Observe all decompositions of $w=xyz$ such that $|x| \leq k$ and $|y| \geq 1$:

\begin{align*}
&x=(aba)^n, \: n \geq 0\\
&y=(aba)^{3m}, \: m\geq1, \: n+m \leq k\\
&z=(aba)^{3k-3m-n}
\end{align*}

Choose $q$ such that $xy^qz \not\in L:$

$$s = xy^qz = (aba)^n(aba)^{3mq}(aba)^{3k-3m-n} = (aba)^{3(k+m(q-1))}$$

Note that $s \in L, \: \forall q\geq0$. This is because regardless of $q$'s value, $aba$ will always be multiplied by a multiple of 3 in the end, due to the presence of 3 on one side of the multiplication of $aba$'s exponent, meaning $aba$'s exponent will always be a multiple of 3. Therefore the Pumping Lemma does not work here, in fact we can derive the regular expression for $L$ since it is regular:

$$L = (abaabaaba)^*$$



% --------------------------------- QUESTION 
\subsection*{1.b}

Assume $L$ is regular, then there exists a pumping length $k$ for $L$. We can choose arbitrary string $s = (01)^k(10)^k, \: |s| > k; s\in L$ since the number of $01$ and $10$ in $s$ are both $k$. 

Observe all decompositions of $s = xyz$ such that $|x| \leq k$ and $|y| \geq 1$:

\begin{align*}
&x=(01)^n, \: n \geq 0\\
&y=(01)^m, \: m\geq1, \: n+m \leq k\\
&z=(01)^{k-m-n}(10)^k
\end{align*}

Choose $q$ such that $xy^qz \not\in L$:

$$s=(01)^n(01)^{mq}(01)^{k-m-n}(10)^k = (01)^{k+m(q-1)}(10)^k$$

If $q=0$, then $s=(01)^{k-m}(10)^k$. $s\in L$ because in every pair of $(10)$ there contains one $(01)$ (i.e. $01 \in 1010$), and in every pair of $(01)$ there exists one $(10)$ (i.e. $10 \in 0101$). 

So if we have $(k-m)$ number of $(01)$ then inside there must exist $(k-m-1)$ number of $(10)$, and the count of $(10)$ is increased to $(k+k-m-1) = 2k-m-1$. 

The reverse is true for $k$ number of $(10)$, since we can count $(k-m-1)$ number of $(10)$ inside the $(k-m)$ number of $(01)$, the count of $(10)$ is also $(k+k-m-1) = 2k -m -1$. Then we see the number of both $(10)$ and $(01)$ are actually the same regardless of their exponents, and the same exists for $q=2$. 

So the pumping lemma fails. In fact, we can prove that $L$ is regular by giving its regular expression. The regular expression is trivial to produce since we know already that $\epsilon, 0, 1$ are already $\in L$ (vacuously true) and our previously testing gives the condition that $w\in L$ if it begins and ends by the same letter:

$$L = \epsilon \cup 1 \cup 0 \cup 0(1\cup0)^*0 \cup 1(1\cup0)^*1$$












\end{document}